\chapter{A File of Python Written as Zeno}\label{ch-convert}

The block of the last chapter keeps code as it was written. Since 0.9
\verb|zeno convert| goes the other way: it takes a whole file of
Python and writes it again as Zeno. What it can say in Zeno it says
in Zeno, and what it cannot it leaves as Python, in an
\verb|inline python| block that the Zeno calls. So what it writes
always runs, and prints what the Python printed.

\begin{note}
The one promise is that what is written is never silently wrong.
Where Zeno would do even a little otherwise than Python, the converter
either says so in a comment or keeps the line as Python.
\end{note}

\section{A first file}

Here is a small till, in Python:

\begin{zeno}[file=shop.py, dir=convert, run=no]
# A till: prices, a basket, and a receipt.
import textwrap

PRICES = {"tea": 2.5, "cake": 3.0, "jam": 1.25}


def line_total(item, count):
    return PRICES[item] * count


def receipt(basket):
    lines = []
    total = 0.0
    for item, count in basket.items():
        cost = line_total(item, count)
        total += cost
        lines.append(f"{item:<6} x{count}  {cost:6.2f}")
    lines.append(f"{'total':<9} {total:6.2f}")
    return lines


basket = {"tea": 2, "jam": 3}
basket["cake"] = 1
for line in receipt(basket):
    print(line)
words = sorted(word.upper() for word in basket)
print(", ".join(words))
note = textwrap.fill("Thank you for shopping with us, and come back soon.", width=24)
print(note)
\end{zeno}

One command writes \verb|shop.zeno| beside it, and says what it did
with each piece of the file:

\begin{zeno}[file=shop.py, dir=convert, show=no, cmd=zeno convert {file}]
# A till: prices, a basket, and a receipt.
import textwrap

PRICES = {"tea": 2.5, "cake": 3.0, "jam": 1.25}


def line_total(item, count):
    return PRICES[item] * count


def receipt(basket):
    lines = []
    total = 0.0
    for item, count in basket.items():
        cost = line_total(item, count)
        total += cost
        lines.append(f"{item:<6} x{count}  {cost:6.2f}")
    lines.append(f"{'total':<9} {total:6.2f}")
    return lines


basket = {"tea": 2, "jam": 3}
basket["cake"] = 1
for line in receipt(basket):
    print(line)
words = sorted(word.upper() for word in basket)
print(", ".join(words))
note = textwrap.fill("Thank you for shopping with us, and come back soon.", width=24)
print(note)
\end{zeno}

This is the Zeno it wrote:

\begin{zeno}[file=read.txt, dir=convert, show=no, cmd=cat shop.zeno]
shop.zeno
\end{zeno}

The types came from the program itself. \verb|line_total| is given a
text and a whole number at its one call, and gives back a float, so it
is \verb|function line_total(item: str, count: int) -> float|.
\verb|basket| holds whole numbers under texts, so \verb|receipt| takes
a \verb|dict[str, int]|. \verb|sorted| is not a function of Zeno, so
the file gets a short one at the top, with a comment that says what it
is. \verb|textwrap| is a library of Python that Zeno has no words for:
its \verb|import| stays in the block, and \verb|fill| is called
through it, as \verb|py.textwrap.fill|.

Both print the same:

\begin{zeno}[file=read.txt, dir=convert, show=no, cmd=zeno run shop.zeno ;; python3 shop.py]
shop.zeno
\end{zeno}

\section{Pieces, and what becomes of each}

The converter cuts the file into \emph{pieces}: the statements that
begin at its left margin. A \verb|def| with its body is one piece, so
is a \verb|class|, a \verb|for| loop at the top of the file, or one
line such as \verb|basket["cake"] = 1|. Each piece ends in one of
three states:

\begin{center}
\begin{tabular}{lP{8cm}}
\toprule
State & What is written \\
\midrule
Zeno & lines of Zeno that mean what the piece meant \\
Zeno with a note & the same, with a \verb|# convert:| comment where
  Zeno does something a reader of Python would not expect \\
Kept & the piece as it was written, in the \verb|inline python|
  block \\
\bottomrule
\end{tabular}
\end{center}

A piece is cut whole, because a piece is what a reader sees as one
thing, and keeping one whole keeps its meaning: nothing inside it is
changed.

\section{What is kept}

Some of Python is not in Zeno: a \verb|lambda|, a function inside a
function, \verb|yield|, a class with a base, \verb|sort(key=...)|. A
piece with one of these is kept, and the report says why, with the
line:

\begin{zeno}[file=counter.py, dir=kept, show=no, cmd=zeno convert {file} ;; cat counter.zeno]
def make_counter():
    count = 0

    def step():
        nonlocal count
        count += 1
        return count

    return step


tick = make_counter()
for _ in range(3):
    tick()
print("ticks:", tick())

values = [3, 1, 2]
values.sort(key=lambda value: -value)
print(values)
print(len(values), max(values))
\end{zeno}

A kept \verb|def| goes at the top of the block. A kept piece that does
something, such as the loop, becomes a function of the block,
\verb|convert_step_2|, and the Zeno calls it where the loop was, so it
runs exactly when it ran in the Python.

Some pieces are kept for what another piece became. \verb|tick| holds
what a kept function gave, which only Python can call, so the line
that gives it its value is kept too: Python can only call Python, and
a variable lives on one side. Each reason in the report says which.

The converter keeps one piece at a time. When several pieces cannot be
written, it keeps the first that has a reason of its own, and tries
the rest again, since a piece that could not be written a moment ago
often can be once the piece it depended on has gone to Python.

\section{What crosses between the two sides}

The Zeno and the block are two sides of the bridge to Python of
\hyperref[ch-other-languages]{Libraries of Other Languages}. A number, a
text, a bool and \verb|none| cross exactly. A value of Python's own,
such as an object a kept class made, stays in Python, and Zeno holds
it as a handle. A list or a dict of Zeno is copied on its way to
Python, and what Python did to its copy would not be seen. So the
converter gives a function of the block only numbers, texts, bools,
\verb|none|, values of Python's own, and lists or dicts written right
there in the call. A line that would give it a list the Zeno holds is
kept.

\section{The notes}

Each note marks a place where the Zeno does what the Python did, but
prints or raises a little otherwise. The comment is written once for
each kind, above the first line it applies to, and the report lists
every line.

\begin{center}
\begin{tabular}{P{4.6cm}P{7.2cm}}
\toprule
Note & What is otherwise \\
\midrule
a value is shown as Zeno shows it & \verb|true|, \verb|false| and
  \verb|none| in small letters, a text inside a list in double quotes,
  a dict between \verb|[| and \verb|]| \\
a very large or very small float & Python's \verb|1e+16| is
  \verb|1e16| \\
the message of an error & in Zeno's words where Zeno raised it \\
an \verb|except| catches more & Zeno raises an error where Python
  raises none: a whole number past 64 bits, say \\
\verb|upper| and \verb|lower| & a to z only \\
\verb|math.round| & a half goes away from zero: 2.5 to 3 \\
a float from the first value & the variable holds a float later, and
  a variable of Zeno has one type \\
\bottomrule
\end{tabular}
\end{center}

\section{What it does not promise}

A program that Python stops with \verb|NameError| is outside the
promise: a name the converter declares before a block has the first
value of its type, where Python would stop. Whole numbers past 64 bits
make Zeno stop with \verb|OverflowError|, so that is never silent, but
\verb|int()| of a float past 64 bits gives the largest whole number
Zeno has. The rest is rare, and is listed in \verb|docs/convert.md|.

\section{How it is checked}

Every rule is held to its word by programs of Python in
\verb|tools/convert| of the repository. Each is converted, run with
\verb|python3| and with \verb|zeno run|, and what the two print is
compared; the first line of each says how many of its pieces are to be
kept, so that a rule that stops working is seen even when the output
is still the same. One of them is a program of 240 lines, with
classes, a dataclass, dicts of objects, exceptions, f-strings and a
report, and all of it but its \verb|import json| is written as Zeno.

Last of all, the converter reads what it wrote with Zeno's own
checker, in the same program, before it writes the file. A piece the
checker does not take is kept, with the checker's words as the reason.
Nothing Zeno would refuse is ever given out.

\begin{note}
C is the next language \verb|zeno convert| is to read, with what it
cannot read left in an \verb|inline c| block.
\end{note}
